\subsection{Analisis Kebutuhan Sistem}
Analisis kebutuhan sistem dirancang untuk memastikan bahwa perangkat lunak memenuhi ekspektasi pengguna dan kriteria fungsional Proyek 2:

\subsubsection{Kebutuhan Fungsional}
\begin{enumerate}
    \item \textbf{SKF-01 (Manajemen Akun):} Sistem wajib menyediakan fitur penambahan, pengubahan, pemilihan akun aktif, serta penghapusan akun portal Wi-Fi.
    \item \textbf{SKF-02 (Enkripsi Kredensial):} Sistem wajib mengenkripsi kata sandi menggunakan AES-256 melalui Android Keystore sebelum disimpan pada media penyimpanan lokal.
    \item \textbf{SKF-03 (Autentikasi Captive Portal):} Sistem wajib mendeteksi keterhubungan ke jaringan Wi-Fi, membaca parameter formulir portal, dan mengirimkan permintaan HTTP POST untuk autentikasi secara otomatis.
    \item \textbf{SKF-04 (Fallback Login Manual):} Sistem wajib menyediakan antarmuka WebView terintegrasi apabila proses otomasi HTTP mengalami kendala (misal: verifikasi Captcha tambahan).
    \item \textbf{SKF-05 (Quick Settings Tile):} Sistem wajib menyediakan tombol pintasan \textit{Tile} di panel kontrol sistem Android untuk login instan.
    \item \textbf{SKF-06 (Autentikasi Biometrik):} Sistem menyediakan opsi penguncian tombol login dengan otorisasi sidik jari atau PIN perangkat.
\end{enumerate}

\subsubsection{Kebutuhan Non-Fungsional}
\begin{enumerate}
    \item \textbf{SKNF-01 (Kinerja):} Proses eksekusi login portal di latar belakang tidak boleh melebihi batas waktu (\textit{timeout}) 5 detik dalam kondisi jaringan stabil.
    \item \textbf{SKNF-02 (Privasi Data):} Kredensial akun tidak boleh dikirimkan ke server pengembang atau pihak ketiga mana pun selain host \textit{captive portal} resmi.
    \item \textbf{SKNF-03 (Kompatibilitas):} Aplikasi wajib dapat dieksekusi dengan baik mulai dari Android versi 10 (API 29) hingga versi 14 (API 34).
\end{enumerate}

\subsection{Arsitektur Sistem}
Aplikasi Coneckin mengadopsi pola arsitektur MVVM (\textit{Model-View-ViewModel}) yang memisahkan tanggung jawab kode menjadi tiga lapisan utama:
\begin{enumerate}
    \item \textbf{View Layer (Jetpack Compose):} Bertanggung jawab menampilkan antarmuka pengguna (\texttt{MainActivity}, \texttt{ManualLoginScreen}, dan komponen dialog). Lapisan ini bersifat reaktif dan mengamati perubahan status data (\textit{state}) tanpa memuat logika bisnis.
    \item \textbf{ViewModel \& Coordinator Layer:} Mengelola siklus hidup status aplikasi, menangani input pengguna, serta mengoordinasikan interaksi antara lapisan antarmuka dan domain bisnis.
    \item \textbf{Data \& Service Layer:} Terdiri dari modul \texttt{AccountStore} (manajemen kredensial terenkripsi), \texttt{CryptoManager} (kriptografi Keystore), \texttt{PortalLogin} (mesin otomasi HTTP handshake), \texttt{LoginTileService} (integrasi panel sistem Android), dan \texttt{Remote} (sinkronisasi direktori portal).
\end{enumerate}

\subsection{Perancangan Alur Otomasi Login}
Alur kerja otomasi autentikasi \textit{captive portal} dirancang dengan algoritma berikut:
\begin{enumerate}
    \item Pengguna menekan tombol \textbf{LOGIN WIFI} pada antarmuka aplikasi atau mengetuk \textit{Quick Settings Tile}.
    \item Aplikasi memeriksa ketersediaan koneksi Wi-Fi menggunakan \texttt{ConnectivityManager}. Jika Wi-Fi belum aktif, proses dibatalkan dengan pesan galat \texttt{NO\_WIFI}.
    \item Aplikasi mengambil kredensial akun aktif yang didekripsi melalui modul \texttt{CryptoManager}.
    \item Aplikasi mengirimkan permintaan deteksi ke URL referensi konektivitas (\textit{probe}). Apabila respons dialihkan (HTTP 302/307) menuju alamat portal (misal \texttt{ypbpi.hotspot}), aplikasi mengekstrak parameter formulir login.
    \item Aplikasi mengirimkan paket HTTP POST berisi nama pengguna dan kata sandi menuju titik akhir autentikasi portal.
    \item Jika respons menghasilkan kode status sukses atau redirection menuju internet publik, aplikasi menandai status \textbf{Berhasil Terhubung} dan memperbarui status \textit{Tile} menjadi aktif.
    \item Jika terjadi kegagalan respons dalam durasi 4 detik, aplikasi menawarkan opsi \textbf{Buka Halaman Portal} melalui WebView.
\end{enumerate}

\subsection{Perancangan Antarmuka Pengguna (UI/UX)}
Antarmuka Coneckin dirancang mengikuti pedoman \textit{Material Design 3} dengan dukungan tema dinamis (Mode Terang dan Mode Gelap). Layar utama menampilkan:
\begin{itemize}
    \item Kartu status koneksi Wi-Fi saat ini (SSID yang terhubung dan status autentikasi).
    \item Tombol utama berukuran besar (\textit{Call-to-Action}) bertuliskan \textbf{LOGIN WIFI}.
    \item Daftar akun yang tersimpan beserta tombol radio untuk memilih akun aktif.
    \item Tombol aksi cepat tambah akun, pengaturan keamanan biometrik, dan riwayat status koneksi.
\end{itemize}
